\documentclass[10pt,a4paper]{amsart}
\usepackage[margin=19mm]{geometry}
\usepackage{amsmath,amssymb,mathtools}
\usepackage[scr=rsfs,cal=euler]{mathalfa}
\usepackage{xfrac}
\usepackage[unicode,hidelinks,bookmarks=false]{hyperref}
\newcommand{\ie}{i.e.\ }
\newcommand{\cf}{cf.\ }
\newcommand{\resp}{respectively\ }
\newcommand{\Subsection}{\subsection}
\providecommand{\nobreakdash}{\nobreak\hbox{-}}
\setlength{\emergencystretch}{2em}
\multlinegap0pt
\usepackage{bm}
\usepackage{bbm}
\usepackage{dsfont}
\usepackage{xspace}
\usepackage{cancel}
\usepackage{mathtools}
\usepackage{amsthm}
\makeatletter
\@ifundefined{thm}{\newtheorem{thm}{Theorem}}{}
\@ifundefined{lem}{\newtheorem{lem}[thm]{Lemma}}{}
\makeatother
\title[Flexible exponents of non-geometric 3-manifolds]{Flexible exponents of non-geometric 3-manifolds}
\author{Jianfeng Lin, Hongbin Sun, Zhongzi Wang}
\date{Source: arXiv:2604.23965v1 (CC BY 4.0)}
\begin{document}
\maketitle

\noindent\emph{Source and licence.} Flexible exponents of non-geometric 3-manifolds. Version arXiv:2604.23965v1 (CC BY 4.0), distributed under the Creative Commons Attribution 4.0 International licence, as recorded in the arXiv licence metadata for this exact version and shown on the abstract page https://arxiv.org/abs/2604.23965v1. Full rights notice: licenses/arxiv-2604.23965-CC-BY-4.0.txt.\par\smallskip

\noindent\textit{Editorial scope.} Counted unit: the Thm whose source label is \texttt{sphericalsummand} in the article (source0.tex lines 763--765, source label \texttt{sphericalsummand}), delivered together with the article's own complete original proof of it (source0.tex lines 767--808, one \texttt{proof} environment of 42 lines). No mathematics is written, repaired or re-derived: statement and proof are the article's own text line for line. Required context retained uncounted: nongeometric (lines 531--538); classify (lines 542--548); cover (lines 558--562); technical (lines 735--742); easy (lines 751--753). Every definition, display and label of those lines is retained unchanged. Omitted from this package and not redistributed: 1--530, 539--541, 549--557, 563--734, 743--750, 754--762, 766--766, 809--1379. Nothing inside a retained excerpt is omitted or altered: every retained body is delivered exactly as the article's lines are written, so no omission receipt of any kind accompanies this package. Citations: the retained excerpts carry no citation command at all, so no bibliography entry of the article is transcribed and no reference is redistributed. Reference adaptation: none was needed and none was made; every reference inside the delivered unit resolves inside it, so no unresolved reference or citation is produced. Printed slips kept literal: no printed word, symbol, hypothesis or number of the counted statement or of its proof was altered. Layout and class: the article's own class is \texttt{amsart} with packages including amssymb, amsfonts, xy, enumerate, amscd, makecell, graphicx, mathrsfs, amsmath, hyperref, booktabs, cleveref, tikz, inputenc; the wrapper is the lane's pinned \texttt{amsart} editorial wrapper at 10pt on A4, which restates the theorem-like environments the retained text needs and adds the article's own macro definitions that the retained text needs (none), with extra wrapper packages (bm, bbm, dsfont, xspace, cancel, mathtools). The delivered numbering is the unit's own. Assets: no retained span contains a figure environment, a graphics-inclusion command, a PDF-page inclusion, a table or a TikZ picture; no image file is copied into this package, the package declares no asset dependency and no image file is redistributed with it. Licence: version arXiv:2604.23965v1 of this article is distributed under the Creative Commons Attribution 4.0 International licence, as recorded in the arXiv licence metadata for this exact version and shown on the abstract page https://arxiv.org/abs/2604.23965v1. A full rights notice is delivered at licenses/arxiv-2604.23965-CC-BY-4.0.txt. Characters: every retained excerpt is pure ASCII, so the delivered engine, XeLaTeX with its default Latin Modern text font, reports no missing character.
\begin{thm}\label{nongeometric}
    Suppose $M$ is a closed, orientable, non-geometric 3-manifold and  \[
        M=(\#^mS^2\times S^1)\#(\#_{i=1}^l P_i),
    \] where each $P_i$ is covered by $S^3$ and $|\pi_1(P_i)|>1$.
  Then $\alpha(M)=2$.
    
    %Then $P_M(L)\lesssim L^2$.
    \end{thm}

\begin{lem}\label{classify}
Any closed, orientable, non-geometric $3$-manifold $$M=(\#^mS^2\times S^1)\#(\#_{i=1}^l P_i)$$ in Theorem \ref{nongeometric} belongs to one of the following two families:
\begin{enumerate}
\item $M=\#^mS^2\times S^1$ with $m\geq 2$,
\item $M=(\#^mS^2\times S^1)\#(\#_{i=1}^l P_i)$ such that $l\geq 1$, $m+l\geq 2$, and $M\ne \mathbb{R}P^3\#\mathbb{R}P^3$.
\end{enumerate}
\end{lem}

\begin{lem}\label{cover}
 Suppose $M$ is a 3-manifold and  \[
        M=(\#^mS^2\times S^1)\#(\#_{i=1}^l P_i).
    \] where each $P_i$ is covered by $S^3$ and $|\pi_1(P_i)|>1$. We also assume that $l\geq 1$, $m+l\geq 2$, and $M\ne \mathbb{R}P^3\# \mathbb{R}P^3$. Then $M$ has a finite regular cover $M'$ that is homeomorphic to $\#^k S^2\times S^1$ with $k\geq 2$.
 \end{lem}

\begin{thm}\label{technical}
Let $M_k=H\cup H'$ be a standard genus-$k$ Heegaard decomposition of $\#^k S^2\times S^1$, and we equip $M_k$ with any Riemannian metric $g$. Then there exists a constant $C>1$ (depending on the metric $g$), such that for any positive integer $n$, there exists a piecewise smooth map $F_n:M_k\to M_k$ that satisfies the following conditions.
\begin{enumerate}
\item $\text{deg}(F_n)=-4n^2-4n=1-(2n+1)^2$.
\item $F_n$ is a $(Cn)$-Lipschitz map.
\item $F_n|_{H'}=id_{H'}$.
\end{enumerate}
\end{thm}

\begin{lem}\label{easy}
For $M=\#^k S^2\times S^1$ with $k\geq 2$, we have $\alpha(M)\geq 2$.
\end{lem}

\begin{thm}\label{sphericalsummand}
Suppose $M=(\#^m S^2\times S^1)\#(\#_{i=1}^l P_i)$ as in Lemma \ref{classify} (2), then $\alpha(M)\geq 2$ holds.
\end{thm}

\begin{proof}
Our $M$ satisfies the assumption of Lemma \ref{cover}, so there exists a finite regular cover $\pi:\tilde{M}\to M$ that is homeomorphic to $\#^kS^2\times S^1$ for some $k\geq 2$. We denote the deck transformation group by $\Gamma$. 

Let $\tilde{M}=H\cup H'$ be one standard genus-$k$ Heegaard decomposition of $\tilde{M}$, then $H$ is a tubular neighborhood of a $3$-valence graph $G\subset H \subset \tilde{M}$. Since $G\subset \tilde{M}$ is a graph in a $3$-manifold, and $\Gamma$ is finite, we can perturb $G$ such that $\gamma\cdot G\cap G=\emptyset$ holds for any $\gamma\in \Gamma\setminus \{e\}$.  So there exists a small neighborhood $H''$ of $G$ in $\tilde{M}$ such that $\gamma\cdot H''\cap H''=\emptyset$ holds for any $\gamma\in \Gamma\setminus \{e\}$, and $H''$ is homeomorphic to the genus-$k$ handlebody. Note that for any $\gamma\in \Gamma$, $\tilde{M}=\gamma\cdot H''\cup (\tilde{M}\setminus \text{int}(\gamma \cdot H''))$ is also a standard genus-$k$ Heegaard decomposition of $\tilde{M}=\#^kS^2\times S^1$.

%By our choice of $H''$, $\pi|_{H''}:H''\to M$ is an embedding, and $\pi(H'')\subset M$ is homeomorphic to the genus-$k$ handlebody. 


%such that its restriction to $\pi(H'')$ is isometric to $H_k$ (defined in the beginning of Section \ref{technicalsection}). This metric $g$ on $M$ pulls back to a $\Gamma$-invariant metric $\tilde{g}$ on $\tilde{M}$. Then the restriction of $\tilde{g}$ to each $\gamma \cdot H''$ is isometric to $H_k$.
We equip $M$ with a Riemannian metric $g$, and it pulls back to a $\Gamma$-invariant metric $\tilde{g}$ on $\tilde{M}$.
By Theorem \ref{technical}, there exists a constant $C>1$, such that for any positive integer $n$, there is a $(Cn)$-Lipschitz map $F_n:\tilde{M}\to \tilde{M}$, satisfying $\text{deg}(F_n)=1-(2n+1)^2$ and $F_n|_{\tilde{M}\setminus H''}=id_{\tilde{M}\setminus H''}$.

Now we use $F_n$ to construct a $\Gamma$-equivariant map $G_n:\tilde{M}\to \tilde{M}$, as 
\begin{equation*}
\begin{aligned}
G_n(x)=\begin{cases}
x & \text{if\ } x\in \tilde{M}\setminus \bigcup_{\gamma\in \Gamma}\text{int}(\gamma \cdot H''),\\
\gamma \cdot F_n(\gamma^{-1}\cdot x) & \text{if\ } x\in \gamma \cdot H'' \text{\ for\ some\ }\gamma\in \Gamma. 
\end{cases}
\end{aligned}
\end{equation*}



%\begin{figure}
 %   \centering
 %   \includegraphics[width=0.7\textwidth]{H_k.pdf}
 %   \caption{Handlebody $H_k$ with $k=5$. Numbered discs decompose $H_k$ to $Y$-shaped solids.}
 %   \label{Figure 1}
%\end{figure}

We need the fact that $\gamma\cdot H''$'s are disjoint from each other to define $G_n$. Since $F_n$ restricts to identity on $\partial H''$, $G_n$ restricts to the identity on $\gamma\cdot \partial H''$, so $G_n$ is well-defined. By the definition of $G_n$, it is routine to check that it is a $\Gamma$-equivariant map, i.e. $G_n(\gamma\cdot x)=\gamma\cdot G_n(x)$ for any $x\in M$ and $\gamma\in \Gamma$. Since $F_n$ is a piecewise smooth $(Cn)$-Lipschitz map and $\Gamma$ acts on $\tilde{M}$ by isometries, $G_n$ is also a $(Cn)$-Lipschitz map.

By Theorem \ref{technical} (1) (3), for any generic point $x_0 \in \tilde{M}\setminus \bigcup_{\gamma\in \Gamma}\gamma \cdot H''$, we have 
$$F_n^{-1}(x_0)=\{x_0,x_1,\cdots,  x_{(2n+1)^2}\},$$ 
such that $x_i\in H''$ for any $i=1,\cdots, (2n+1)^2$. Moreover, the local degree of $F_n$ at $x_0$ is $1$, and the local degree of $F_n$ at $x_i$ is $-1$ for any $i=1,\cdots, (2n+1)^2$. Then for any generic $x_0\in \tilde{M}\setminus \bigcup_{\gamma\in \Gamma}\gamma \cdot H''$, its $\Gamma$-orbit is contained in $\tilde{M}\setminus \bigcup_{\gamma\in \Gamma}\gamma \cdot H''$ and
$$G_n^{-1}(x_0)=\{x_0\}\cup \Big(\cup_{\gamma\in \Gamma}\gamma\cdot (F_n^{-1}(\gamma^{-1}x_0)\cap H''))\Big).$$
The local degree of $G_n$ at $x_0$ is $1$. Each $\gamma\cdot (F_n^{-1}(\gamma^{-1}x_0)\cap H'')$ has cardinality $(2n+1)^2$, and the local degree of $G_n$ at each such point is $-1$. So we have 
$$\text{deg}(G_n)=1-(2n+1)^2\cdot |\Gamma|.$$

Since $G_n:\tilde{M} \to \tilde{M}$ is $\Gamma$-equivariant, it decends to a map $K_n:M\to M$, such that $\text{deg}(K_n)=\text{deg}(G_n)=1-(2n+1)^2\cdot |\Gamma|$. Since the metric $\tilde{g}$ on $\tilde{M}$ is the pull-back of the metric $g$ on $M$, $K_n$ is also a piecewise smooth $(Cn)$-Lipschitz map.  These facts on the sequence of maps $K_n: M\to M$ imply that $\alpha(M)\geq 2$ holds, as in the proof of Lemma \ref{easy}.
\end{proof}

\end{document}
